\documentclass[11pt,a4paper]{article}
\usepackage[top=2cm, bottom=2.3cm, left=1.8cm, right=1.8cm, headsep=0.8cm]{geometry}
\usepackage{amsmath,amssymb,amsfonts}
\usepackage{enumitem}
\usepackage{fancyhdr}
\usepackage[table]{xcolor}
\usepackage{tcolorbox}
\tcbuselibrary{breakable,skins}
\usepackage{tabularx}
\usepackage{multicol}

% High-Contrast Modern Color Palette
\definecolor{primaryDark}{RGB}{12, 32, 84}       % Midnight Navy
\definecolor{primaryBlue}{RGB}{18, 72, 160}      % Solid Royal Blue
\definecolor{goldAccent}{RGB}{255, 215, 0}       % Gold highlight
\definecolor{l1Green}{RGB}{14, 120, 62}          % Vivid Forest Green
\definecolor{l2Amber}{RGB}{205, 80, 0}           % Rich Amber
\definecolor{cardBg}{RGB}{250, 252, 255}         % Crisp Card Background
\definecolor{cardBorder}{RGB}{130, 170, 220}     % Steel Blue Border
\definecolor{tableStripe}{RGB}{242, 246, 252}    % Alternating Table Row
\definecolor{tableDarkHeader}{RGB}{12, 32, 84}   % Dark Table Header

% Page styling
\pagestyle{fancy}
\fancyhf{}
\setlength{\headheight}{16pt}
\fancyhead[L]{\small\textsf{\textbf{\textcolor{primaryDark}{Sohan Mocks Series}} $\vert$ \textbf{\textcolor{primaryBlue}{Mathematics Full Mock}}}}
\fancyhead[R]{\small\textsf{\textbf{\textcolor{l1Green}{Max Marks: 100 $\vert$ 60 Mins}}}}
\fancyfoot[C]{\textbf{\textsf{\thepage}}}
\renewcommand{\headrulewidth}{0.8pt}
\renewcommand{\footrulewidth}{0.4pt}

% Custom tcolorbox styles
\tcbset{
  dayheader/.style={
    colback=primaryDark,
    coltext=white,
    boxrule=0pt,
    arc=3mm,
    left=14pt, right=14pt, top=10pt, bottom=10pt
  },
  sectionbanner/.style={
    colback=primaryBlue,
    coltext=white,
    boxrule=0pt,
    arc=2mm,
    left=10pt, right=10pt, top=6pt, bottom=6pt
  },
  solbox/.style={
    enhanced,
    breakable,
    colback=cardBg,
    colframe=cardBorder,
    boxrule=1.0pt,
    arc=2mm,
    left=10pt, right=10pt, top=7pt, bottom=7pt
  }
}

\begin{document}

% ==================== TEST HEADER BANNER ====================
\begin{tcolorbox}[dayheader]
\begin{center}
  {\Large \textbf{\textsf{\textcolor{goldAccent}{SOHAN MOCKS TEST SERIES}}}}\\[3pt]
  {\LARGE \textbf{\textsf{FULL SYLLABUS: MATHEMATICS (JEE MAIN MOCK TEST)}}}\\[5pt]
  {\normalsize \textbf{\textsf{Domains:}} Calculus $\vert$ Algebra $\vert$ Coordinate Geometry $\vert$ Vectors \& 3D $\vert$ Trigonometry}\\[3pt]
  {\footnotesize \textbf{Exam Pattern:} 25 Questions (20 MCQs $+4/-1$ + 5 Numerical $+4/0$) \quad $\vert$ \quad \textbf{Time:} 60 Minutes \quad $\vert$ \quad \textbf{Max Marks:} 100}
\end{center}
\end{tcolorbox}
\vspace{6pt}

% ==================== SECTION I: SINGLE CORRECT CHOICE ====================
\begin{tcolorbox}[sectionbanner]
  {\large \textbf{\textsf{SECTION I: SINGLE CORRECT CHOICE TYPE (Q01 -- Q20)}}} \hfill {\small \textbf{Marking: $+4$ / $-1$}}
\end{tcolorbox}
\vspace{4pt}
\begin{enumerate}[label=\textbf{\arabic*.},leftmargin=18pt,itemsep=12pt]

  \item Let $f: \mathbb{R} \to (1, 3)$ be a bijective function defined by $f(x) = \dfrac{e^x - e^{-x}}{e^x + e^{-x}} + 2$. If $f^{-1}(x)$ denotes the inverse function of $f(x)$, then $f^{-1}(x)$ is equal to:
  \begin{multicols}{2}
    \begin{enumerate}[label=(\Alph*)]
      \item $\dfrac{1}{2}\ln\left(\dfrac{x - 1}{3 - x}\right)$
      \item $\dfrac{1}{2}\ln\left(\dfrac{3 - x}{x - 1}\right)$
      \item $\ln\left(\dfrac{x - 1}{3 - x}\right)$
      \item $\dfrac{1}{2}\ln\left(\dfrac{x + 1}{3 + x}\right)$
    \end{enumerate}
  \end{multicols}

  \item Let $A$ be a $3 \times 3$ non-singular matrix with real entries such that $\det(A) = 3$. If matrix $B = \operatorname{adj}(2A)$, then the value of $\det(\operatorname{adj}(B))$ is:
  \begin{multicols}{2}
    \begin{enumerate}[label=(\Alph*)]
      \item $2^8 \times 3^4$
      \item $2^{12} \times 3^4$
      \item $2^{10} \times 3^3$
      \item $2^{12} \times 3^2$
    \end{enumerate}
  \end{multicols}

  \item If a complex number $z = x + iy$ ($x, y \in \mathbb{R}$) satisfies the equation $\left|\dfrac{z - 2i}{z + 2i}\right| = \sqrt{2}$, then the locus of $z$ in the Argand plane is a circle whose center and radius are respectively:
  \begin{multicols}{2}
    \begin{enumerate}[label=(\Alph*)]
      \item $(0, -6)$ and $4\sqrt{2}$
      \item $(0, 6)$ and $4\sqrt{2}$
      \item $(0, -6)$ and $2\sqrt{7}$
      \item $(0, -4)$ and $2\sqrt{3}$
    \end{enumerate}
  \end{multicols}

  \item The value of the limit $\lim_{x \to 0} \dfrac{e^{x^2} - \cos x - \dfrac{3}{2}x^2}{x^4}$ is equal to:
  \begin{multicols}{2}
    \begin{enumerate}[label=(\Alph*)]
      \item $\dfrac{11}{24}$
      \item $\dfrac{5}{12}$
      \item $\dfrac{7}{24}$
      \item $\dfrac{1}{2}$
    \end{enumerate}
  \end{multicols}

  \item Let $\alpha$ and $\beta$ be the roots of the quadratic equation $x^2 - 6x - 2 = 0$ with $\alpha > \beta$. If $a_n = \alpha^n - \beta^n$ for all integers $n \ge 1$, then the value of $\dfrac{a_{10} - 2a_8}{2a_9}$ is:
  \begin{multicols}{2}
    \begin{enumerate}[label=(\Alph*)]
      \item $1$
      \item $2$
      \item $3$
      \item $6$
    \end{enumerate}
  \end{multicols}

  \item The value of the definite integral $I = \displaystyle\int_{0}^{\pi} \dfrac{x \sin x}{1 + \cos^2 x}\, dx$ is:
  \begin{multicols}{2}
    \begin{enumerate}[label=(\Alph*)]
      \item $\dfrac{\pi^2}{4}$
      \item $\dfrac{\pi^2}{2}$
      \item $\dfrac{\pi^2}{8}$
      \item $\pi^2$
    \end{enumerate}
  \end{multicols}

  \item The sum of the infinite series $S = \displaystyle\sum_{n=1}^{\infty} \dfrac{2n + 1}{n^2(n + 1)^2}$ is equal to:
  \begin{multicols}{2}
    \begin{enumerate}[label=(\Alph*)]
      \item $\dfrac{1}{2}$
      \item $1$
      \item $\dfrac{3}{2}$
      \item $2$
    \end{enumerate}
  \end{multicols}

  \item If the straight lines $L_1: \dfrac{x - 1}{2} = \dfrac{y + 1}{3} = \dfrac{z - 1}{4}$ and $L_2: \dfrac{x - 3}{1} = \dfrac{y - k}{2} = \dfrac{z}{1}$ intersect at a point, then the value of $k$ is:
  \begin{multicols}{2}
    \begin{enumerate}[label=(\Alph*)]
      \item $\dfrac{7}{2}$
      \item $\dfrac{9}{2}$
      \item $5$
      \item $\dfrac{11}{2}$
    \end{enumerate}
  \end{multicols}

  \item The solution curve of the differential equation $\dfrac{dy}{dx} + y \tan x = 2x + x^2 \tan x$ passing through $(0, 1)$ satisfies $y\left(\dfrac{\pi}{4}\right)$ equal to:
  \begin{multicols}{2}
    \begin{enumerate}[label=(\Alph*)]
      \item $\dfrac{\pi^2}{16} + \dfrac{1}{\sqrt{2}}$
      \item $\dfrac{\pi^2}{16} - \dfrac{1}{\sqrt{2}}$
      \item $\dfrac{\pi^2}{8} + \dfrac{1}{\sqrt{2}}$
      \item $\dfrac{\pi^2}{16} + 1$
    \end{enumerate}
  \end{multicols}

  \item The mean and variance of 7 observations are $8$ and $16$, respectively. If five of the observations are $2, 4, 10, 12, 14$, then the absolute difference between the remaining two observations is:
  \begin{multicols}{2}
    \begin{enumerate}[label=(\Alph*)]
      \item $1$
      \item $2$
      \item $3$
      \item $4$
    \end{enumerate}
  \end{multicols}

  \item If a line $y = mx + c$ with slope $m > 0$ is a common tangent to both the parabola $y^2 = 8x$ and the circle $x^2 + y^2 = 2$, then the value of $m + c$ is:
  \begin{multicols}{2}
    \begin{enumerate}[label=(\Alph*)]
      \item $1$
      \item $2$
      \item $3$
      \item $4$
    \end{enumerate}
  \end{multicols}

  \item The value of the finite sum $\displaystyle\sum_{r=0}^{20} (-1)^r \, \dfrac{^{20}C_r}{r + 1}$ is equal to:
  \begin{multicols}{2}
    \begin{enumerate}[label=(\Alph*)]
      \item $\dfrac{1}{20}$
      \item $\dfrac{1}{21}$
      \item $\dfrac{1}{22}$
      \item $0$
    \end{enumerate}
  \end{multicols}

  \item Let $\vec{a} = \hat{i} + \hat{j} + \hat{k}$, $\vec{b} = \hat{i} - \hat{j} + 2\hat{k}$, and $\vec{c} = x\hat{i} + (x - 2)\hat{j} - \hat{k}$. If $\vec{c}$ lies in the plane of $\vec{a}$ and $\vec{b}$, then the value of $x$ is:
  \begin{multicols}{2}
    \begin{enumerate}[label=(\Alph*)]
      \item $-2$
      \item $2$
      \item $-1$
      \item $0$
    \end{enumerate}
  \end{multicols}

  \item The area (in square units) of the region bounded by the parabola $y = x^2$ and the straight line $y = 2x + 3$ is:
  \begin{multicols}{2}
    \begin{enumerate}[label=(\Alph*)]
      \item $\dfrac{16}{3}$
      \item $\dfrac{32}{3}$
      \item $\dfrac{28}{3}$
      \item $10$
    \end{enumerate}
  \end{multicols}

  \item The total number of 5-digit numbers that can be formed using the digits $\{1, 2, 3, 4, 5, 6, 7\}$ without repetition such that the number is divisible by $4$ is:
  \begin{multicols}{2}
    \begin{enumerate}[label=(\Alph*)]
      \item $480$
      \item $600$
      \item $720$
      \item $840$
    \end{enumerate}
  \end{multicols}

  \item The maximum area of a rectangle inscribed in the ellipse $\dfrac{x^2}{16} + \dfrac{y^2}{9} = 1$ with sides parallel to the coordinate axes is:
  \begin{multicols}{2}
    \begin{enumerate}[label=(\Alph*)]
      \item $12$
      \item $24$
      \item $36$
      \item $48$
    \end{enumerate}
  \end{multicols}

  \item Let $f(x) = \begin{cases} x^2 \sin\left(\dfrac{1}{x}\right) + ax, & x \ne 0 \\ b, & x = 0 \end{cases}$. If $f(x)$ is differentiable at $x = 0$ with $f'(0) = 3$, then the ordered pair $(a, b)$ is:
  \begin{multicols}{2}
    \begin{enumerate}[label=(\Alph*)]
      \item $(3, 0)$
      \item $(0, 3)$
      \item $(3, 1)$
      \item $(1, 0)$
    \end{enumerate}
  \end{multicols}

  \item The value of $\tan\left( \displaystyle\sum_{r=1}^{10} \tan^{-1}\left(\dfrac{1}{r^2 + r + 1}\right) \right)$ is equal to:
  \begin{multicols}{2}
    \begin{enumerate}[label=(\Alph*)]
      \item $\dfrac{5}{6}$
      \item $\dfrac{10}{11}$
      \item $\dfrac{11}{12}$
      \item $\dfrac{6}{7}$
    \end{enumerate}
  \end{multicols}

  \item The length of the tangent drawn from any point on the circle $C_1: x^2 + y^2 + 4x + 6y - 12 = 0$ to the circle $C_2: x^2 + y^2 + 4x + 6y - 4 = 0$ is:
  \begin{multicols}{2}
    \begin{enumerate}[label=(\Alph*)]
      \item $2$
      \item $2\sqrt{2}$
      \item $3$
      \item $4$
    \end{enumerate}
  \end{multicols}

  \item Box I contains $4$ red and $6$ black balls, while Box II contains $5$ red and $5$ black balls. A box is chosen uniformly at random and a ball is drawn. If the drawn ball is found to be red, what is the probability that it originated from Box I?
  \begin{multicols}{2}
    \begin{enumerate}[label=(\Alph*)]
      \item $\dfrac{4}{9}$
      \item $\dfrac{5}{9}$
      \item $\dfrac{2}{5}$
      \item $\dfrac{1}{2}$
    \end{enumerate}
  \end{multicols}
\end{enumerate}
\vspace{10pt}

% ==================== SECTION II: NUMERICAL VALUE TYPE ====================
\begin{tcolorbox}[sectionbanner]
  {\large \textbf{\textsf{SECTION II: NUMERICAL VALUE TYPE (Q21 -- Q25)}}} \hfill {\small \textbf{Marking: $+4$ / $0$}}
\end{tcolorbox}
\vspace{4pt}
\begin{enumerate}[label=\textbf{\arabic*.},leftmargin=18pt,itemsep=12pt,start=21]

  \item If a differentiable function $f(x)$ satisfies the integral equation $\displaystyle\int_{0}^{x^2} f(t)\, dt = \dfrac{2}{5}x^5 + \dfrac{4}{3}x^3$ for all $x > 0$, then find the value of $f(4)$.

  \item Let $A = \begin{bmatrix} 1 & 2 \\ 0 & 1 \end{bmatrix}$. If $A^{10} = \begin{bmatrix} 1 & k \\ 0 & 1 \end{bmatrix}$, then the value of $k$ is:

  \item Let $\vec{a} = 2\hat{i} + \hat{j} - 2\hat{k}$ and $\vec{b} = \hat{i} + \hat{j}$. If $\vec{c}$ is a vector such that $\vec{a} \cdot \vec{c} = |\vec{c}|$, $|\vec{c} - \vec{a}| = 2\sqrt{2}$, and $\vec{c}$ is perpendicular to $\vec{a} \times \vec{b}$, then find the value of $|(\vec{a} \times \vec{b}) \times \vec{c}|$.

  \item A focal chord of the parabola $y^2 = 4x$ passing through its focus $S(1, 0)$ intersects the curve at points $P(x_1, y_1)$ and $Q(x_2, y_2)$. If the length of the focal chord $PQ$ is $16$, then find the value of $x_1 + x_2$.

  \item If the roots of the cubic equation $x^3 - 9x^2 + 23x - c = 0$ are three positive integers forming an arithmetic progression, find the value of the constant $c$.
\end{enumerate}

\newpage

% ==================== MASTER ANSWER KEY ====================
\begin{tcolorbox}[sectionbanner]
  {\large \textbf{\textsf{MASTER 100\% VERIFIED ANSWER KEY}}}
\end{tcolorbox}
\vspace{4pt}

\begin{center}
\renewcommand{\arraystretch}{1.25}
\small
\begin{tabularx}{\textwidth}{|>{\centering\arraybackslash}c|>{\centering\arraybackslash}c|>{\centering\arraybackslash}c|>{\centering\arraybackslash}c|>{\centering\arraybackslash}c|>{\centering\arraybackslash}c|>{\centering\arraybackslash}c|>{\centering\arraybackslash}c|>{\centering\arraybackslash}c|>{\centering\arraybackslash}X|}
\hline
\rowcolor{tableDarkHeader}
\textcolor{white}{\textbf{Q.}} & \textcolor{white}{\textbf{Ans}} & \textcolor{white}{\textbf{Q.}} & \textcolor{white}{\textbf{Ans}} & \textcolor{white}{\textbf{Q.}} & \textcolor{white}{\textbf{Ans}} & \textcolor{white}{\textbf{Q.}} & \textcolor{white}{\textbf{Ans}} & \textcolor{white}{\textbf{Q.}} & \textcolor{white}{\textbf{Ans}} \\
\hline
\rowcolor{white}
\textbf{01} & A & \textbf{06} & A & \textbf{11} & C & \textbf{16} & B & \textbf{21} & \textbf{12} \\
\hline
\rowcolor{tableStripe}
\textbf{02} & B & \textbf{07} & B & \textbf{12} & B & \textbf{17} & A & \textbf{22} & \textbf{20} \\
\hline
\rowcolor{white}
\textbf{03} & A & \textbf{08} & B & \textbf{13} & A & \textbf{18} & A & \textbf{23} & \textbf{3} \\
\hline
\rowcolor{tableStripe}
\textbf{04} & A & \textbf{09} & A & \textbf{14} & B & \textbf{19} & B & \textbf{24} & \textbf{14} \\
\hline
\rowcolor{white}
\textbf{05} & C & \textbf{10} & B & \textbf{15} & B & \textbf{20} & A & \textbf{25} & \textbf{15} \\
\hline
\end{tabularx}
\end{center}

\vspace{10pt}

% ==================== HINTS \& STEP-BY-STEP SOLUTIONS ====================
\begin{tcolorbox}[sectionbanner]
  {\large \textbf{\textsf{HINTS \& STEP-BY-STEP DETAILED SOLUTIONS}}}
\end{tcolorbox}
\vspace{6pt}

\begin{tcolorbox}[solbox]
\textbf{\textcolor{primaryDark}{Question 01:}} \hfill \textbf{\textcolor{l1Green}{Correct Answer: A}}\\
\vspace{4pt}
Let $y = f(x) = \dfrac{e^x - e^{-x}}{e^x + e^{-x}} + 2 = \dfrac{e^{2x} - 1}{e^{2x} + 1} + 2$.

Subtracting 2 from both sides:
$$y - 2 = \frac{e^{2x} - 1}{e^{2x} + 1}$$

Applying componendo and dividendo (or cross-multiplication):
$$(y - 2)(e^{2x} + 1) = e^{2x} - 1$$
$$e^{2x}(y - 2 - 1) = -1 - (y - 2)$$
$$e^{2x}(y - 3) = -(y - 1) = 1 - y$$
$$e^{2x} = \frac{1 - y}{y - 3} = \frac{y - 1}{3 - y}$$

Taking natural logarithm of both sides for $y \in (1, 3)$:
$$2x = \ln\left(\frac{y - 1}{3 - y}\right) \implies x = \frac{1}{2}\ln\left(\frac{y - 1}{3 - y}\right)$$

Therefore, the inverse function is:
$$f^{-1}(x) = \frac{1}{2}\ln\left(\frac{x - 1}{3 - x}\right)$$

Hence, the correct option is \textbf{(A)}.
\end{tcolorbox}

\begin{tcolorbox}[solbox]
\textbf{\textcolor{primaryDark}{Question 02:}} \hfill \textbf{\textcolor{l1Green}{Correct Answer: B}}\\
\vspace{4pt}
For any $n \times n$ matrix $M$, we have the standard determinant properties:
1. $\det(kM) = k^n \det(M)$
2. $\det(\operatorname{adj}(M)) = (\det(M))^{n-1}$

Here $n = 3$, so $\det(\operatorname{adj}(M)) = (\det(M))^2$.

First, evaluate $\det(2A)$:
$$\det(2A) = 2^3 \det(A) = 8 \times 3 = 24$$

Now, for $B = \operatorname{adj}(2A)$:
$$\det(B) = \det(\operatorname{adj}(2A)) = (\det(2A))^{3-1} = (24)^2$$

Next, evaluate $\det(\operatorname{adj}(B))$:
$$\det(\operatorname{adj}(B)) = (\det(B))^{3-1} = (\det(B))^2 = \left(24^2\right)^2 = 24^4$$

Decomposing 24 into prime factors:
$$24 = 2^3 \times 3$$
$$24^4 = (2^3 \times 3)^4 = 2^{12} \times 3^4$$

Hence, the correct option is \textbf{(B)}.
\end{tcolorbox}

\begin{tcolorbox}[solbox]
\textbf{\textcolor{primaryDark}{Question 03:}} \hfill \textbf{\textcolor{l1Green}{Correct Answer: A}}\\
\vspace{4pt}
Given the equation:
$$\left|\frac{z - 2i}{z + 2i}\right| = \sqrt{2} \implies |z - 2i| = \sqrt{2}|z + 2i|$$

Squaring both sides:
$$|z - 2i|^2 = 2|z + 2i|^2$$

Substitute $z = x + iy$:
$$|x + i(y - 2)|^2 = 2|x + i(y + 2)|^2$$
$$x^2 + (y - 2)^2 = 2\left[x^2 + (y + 2)^2\right]$$
$$x^2 + y^2 - 4y + 4 = 2x^2 + 2y^2 + 8y + 8$$

Rearranging terms to one side:
$$x^2 + y^2 + 12y + 4 = 0$$

Comparing with the standard circle equation $x^2 + y^2 + 2gx + 2fy + c = 0$:
$$g = 0, \quad f = 6, \quad c = 4$$

- \textbf{Center}: $(-g, -f) = (0, -6)$
- \textbf{Radius}: $r = \sqrt{g^2 + f^2 - c} = \sqrt{0^2 + 6^2 - 4} = \sqrt{36 - 4} = \sqrt{32} = 4\sqrt{2}$

Hence, the correct option is \textbf{(A)}.
\end{tcolorbox}

\begin{tcolorbox}[solbox]
\textbf{\textcolor{primaryDark}{Question 04:}} \hfill \textbf{\textcolor{l1Green}{Correct Answer: A}}\\
\vspace{4pt}
Using Maclaurin series expansions about $x = 0$:
$$e^{u} = 1 + u + \frac{u^2}{2!} + O(u^3) \implies e^{x^2} = 1 + x^2 + \frac{x^4}{2} + O(x^6)$$
$$\cos x = 1 - \frac{x^2}{2!} + \frac{x^4}{4!} + O(x^6) = 1 - \frac{x^2}{2} + \frac{x^4}{24} + O(x^6)$$

Substituting these expansions into the numerator:
$$\text{Numerator} = \left(1 + x^2 + \frac{x^4}{2}\right) - \left(1 - \frac{x^2}{2} + \frac{x^4}{24}\right) - \frac{3}{2}x^2$$
$$= 1 - 1 + \left(x^2 + \frac{x^2}{2} - \frac{3}{2}x^2\right) + \left(\frac{1}{2} - \frac{1}{24}\right)x^4 + O(x^6)$$
$$= 0 + 0 + \left(\frac{12 - 1}{24}\right)x^4 + O(x^6) = \frac{11}{24}x^4 + O(x^6)$$

Dividing by the denominator $x^4$:
$$\lim_{x \to 0} \frac{\dfrac{11}{24}x^4 + O(x^6)}{x^4} = \frac{11}{24}$$

Hence, the correct option is \textbf{(A)}.
\end{tcolorbox}

\begin{tcolorbox}[solbox]
\textbf{\textcolor{primaryDark}{Question 05:}} \hfill \textbf{\textcolor{l1Green}{Correct Answer: C}}\\
\vspace{4pt}
Since $\alpha$ and $\beta$ are the roots of $x^2 - 6x - 2 = 0$, they satisfy the equation:
$$\alpha^2 - 6\alpha - 2 = 0 \implies \alpha^2 - 2 = 6\alpha$$
$$\beta^2 - 6\beta - 2 = 0 \implies \beta^2 - 2 = 6\beta$$

Multiplying the first equation by $\alpha^8$ and the second by $\beta^8$:
$$\alpha^{10} - 2\alpha^8 = 6\alpha^9$$
$$\beta^{10} - 2\beta^8 = 6\beta^9$$

Subtracting the second equation from the first:
$$(\alpha^{10} - \beta^{10}) - 2(\alpha^8 - \beta^8) = 6(\alpha^9 - \beta^9)$$

By definition of $a_n = \alpha^n - \beta^n$:
$$a_{10} - 2a_8 = 6a_9$$

Dividing both sides by $2a_9$:
$$\frac{a_{10} - 2a_8}{2a_9} = \frac{6a_9}{2a_9} = 3$$

Hence, the correct option is \textbf{(C)}.
\end{tcolorbox}

\begin{tcolorbox}[solbox]
\textbf{\textcolor{primaryDark}{Question 06:}} \hfill \textbf{\textcolor{l1Green}{Correct Answer: A}}\\
\vspace{4pt}
Given the integral:
$$I = \int_{0}^{\pi} \frac{x \sin x}{1 + \cos^2 x}\, dx \quad \text{--- (1)}$$

Applying King's property $\displaystyle\int_a^b f(x)\, dx = \int_a^b f(a + b - x)\, dx$ with $x \to \pi - x$:
$$I = \int_{0}^{\pi} \frac{(\pi - x) \sin(\pi - x)}{1 + \cos^2(\pi - x)}\, dx$$
Since $\sin(\pi - x) = \sin x$ and $\cos(\pi - x) = -\cos x \implies \cos^2(\pi - x) = \cos^2 x$:
$$I = \int_{0}^{\pi} \frac{(\pi - x) \sin x}{1 + \cos^2 x}\, dx = \pi \int_{0}^{\pi} \frac{\sin x}{1 + \cos^2 x}\, dx - I \quad \text{--- (2)}$$

Adding equations (1) and (2):
$$2I = \pi \int_{0}^{\pi} \frac{\sin x}{1 + \cos^2 x}\, dx$$

Substitute $u = \cos x \implies du = -\sin x\, dx$:
When $x = 0$, $u = 1$; when $x = \pi$, $u = -1$.
$$2I = \pi \int_{1}^{-1} \frac{-du}{1 + u^2} = \pi \int_{-1}^{1} \frac{du}{1 + u^2} = 2\pi \int_{0}^{1} \frac{du}{1 + u^2}$$
$$2I = 2\pi \left[ \tan^{-1}(u) \right]_{0}^{1} = 2\pi \left( \frac{\pi}{4} - 0 \right) = \frac{\pi^2}{2}$$
$$I = \frac{\pi^2}{4}$$

Hence, the correct option is \textbf{(A)}.
\end{tcolorbox}

\begin{tcolorbox}[solbox]
\textbf{\textcolor{primaryDark}{Question 07:}} \hfill \textbf{\textcolor{l1Green}{Correct Answer: B}}\\
\vspace{4pt}
Express the numerator in terms of the factors in the denominator:
$$2n + 1 = (n + 1)^2 - n^2 = (n^2 + 2n + 1) - n^2 = 2n + 1$$

Thus, the general term $T_n$ can be decomposed into partial fractions:
$$T_n = \frac{(n + 1)^2 - n^2}{n^2(n + 1)^2} = \frac{(n + 1)^2}{n^2(n + 1)^2} - \frac{n^2}{n^2(n + 1)^2} = \frac{1}{n^2} - \frac{1}{(n + 1)^2}$$

The sum of the first $k$ terms telescopes:
$$S_k = \sum_{n=1}^k \left( \frac{1}{n^2} - \frac{1}{(n + 1)^2} \right)$$
$$S_k = \left(1 - \frac{1}{4}\right) + \left(\frac{1}{4} - \frac{1}{9}\right) + \dots + \left(\frac{1}{k^2} - \frac{1}{(k+1)^2}\right)$$
$$S_k = 1 - \frac{1}{(k + 1)^2}$$

Taking the limit as $k \to \infty$:
$$S = \lim_{k \to \infty} \left( 1 - \frac{1}{(k + 1)^2} \right) = 1 - 0 = 1$$

Hence, the correct option is \textbf{(B)}.
\end{tcolorbox}

\begin{tcolorbox}[solbox]
\textbf{\textcolor{primaryDark}{Question 08:}} \hfill \textbf{\textcolor{l1Green}{Correct Answer: B}}\\
\vspace{4pt}
Two lines intersect if and only if they are coplanar. The condition for coplanarity of two lines with points $\vec{a}_1, \vec{a}_2$ and direction vectors $\vec{b}_1, \vec{b}_2$ is:
$$\begin{vmatrix} x_2 - x_1 & y_2 - y_1 & z_2 - z_1 \\ l_1 & m_1 & n_1 \\ l_2 & m_2 & n_2 \end{vmatrix} = 0$$

From the given equations:
- Point on $L_1$: $(x_1, y_1, z_1) = (1, -1, 1)$, Direction ratios: $\langle 2, 3, 4 \rangle$
- Point on $L_2$: $(x_2, y_2, z_2) = (3, k, 0)$, Direction ratios: $\langle 1, 2, 1 \rangle$

Connecting vector: $(x_2 - x_1, y_2 - y_1, z_2 - z_1) = (3 - 1, k - (-1), 0 - 1) = (2, k + 1, -1)$.

Setting the coplanarity determinant to zero:
$$\begin{vmatrix} 2 & k + 1 & -1 \\ 2 & 3 & 4 \\ 1 & 2 & 1 \end{vmatrix} = 0$$

Expanding along the first row:
$$2 [3(1) - 4(2)] - (k + 1)[2(1) - 4(1)] + (-1)[2(2) - 3(1)] = 0$$
$$2(3 - 8) - (k + 1)(2 - 4) - 1(4 - 3) = 0$$
$$2(-5) - (k + 1)(-2) - 1(1) = 0$$
$$-10 + 2k + 2 - 1 = 0$$
$$2k - 9 = 0 \implies 2k = 9 \implies k = \frac{9}{2}$$

Hence, the correct option is \textbf{(B)}.
\end{tcolorbox}

\begin{tcolorbox}[solbox]
\textbf{\textcolor{primaryDark}{Question 09:}} \hfill \textbf{\textcolor{l1Green}{Correct Answer: A}}\\
\vspace{4pt}
The given differential equation is a first-order linear differential equation of the form:
$$\frac{dy}{dx} + P(x) y = Q(x)$$
where $P(x) = \tan x$ and $Q(x) = 2x + x^2 \tan x$.

1. \textbf{Integrating Factor (I.F.)}:
$$\text{I.F.} = e^{\int P(x)\, dx} = e^{\int \tan x\, dx} = e^{\ln|\sec x|} = \sec x$$

2. \textbf{General Solution}:
$$y \cdot \text{I.F.} = \int Q(x) \cdot \text{I.F.}\, dx + C$$
$$y \sec x = \int (2x + x^2 \tan x)\sec x\, dx + C = \int (2x \sec x + x^2 \sec x \tan x)\, dx + C$$

Notice that $\dfrac{d}{dx}\left(x^2 \sec x\right) = 2x \sec x + x^2 \sec x \tan x$. Hence:
$$y \sec x = x^2 \sec x + C \implies y = x^2 + C \cos x$$

3. \textbf{Initial Condition}:
Given $y(0) = 1$:
$$1 = 0^2 + C \cos(0) = C \implies C = 1$$
Therefore, the particular solution is:
$$y(x) = x^2 + \cos x$$

4. \textbf{Evaluate at $x = \pi/4$}:
$$y\left(\frac{\pi}{4}\right) = \left(\frac{\pi}{4}\right)^2 + \cos\left(\frac{\pi}{4}\right) = \frac{\pi^2}{16} + \frac{1}{\sqrt{2}}$$

Hence, the correct option is \textbf{(A)}.
\end{tcolorbox}

\begin{tcolorbox}[solbox]
\textbf{\textcolor{primaryDark}{Question 10:}} \hfill \textbf{\textcolor{l1Green}{Correct Answer: B}}\\
\vspace{4pt}
Let the two unknown observations be $a$ and $b$.

1. \textbf{Mean condition}:
$$\bar{x} = \frac{2 + 4 + 10 + 12 + 14 + a + b}{7} = 8$$
$$42 + a + b = 56 \implies a + b = 14 \quad \text{--- (1)}$$

2. \textbf{Variance condition}:
$$\sigma^2 = \frac{\sum x_i^2}{N} - (\bar{x})^2 = 16$$
$$\frac{\sum x_i^2}{7} - 8^2 = 16 \implies \frac{\sum x_i^2}{7} = 16 + 64 = 80 \implies \sum x_i^2 = 560$$

Sum of squares of the known 5 observations:
$$2^2 + 4^2 + 10^2 + 12^2 + 14^2 = 4 + 16 + 100 + 144 + 196 = 460$$
Therefore:
$$a^2 + b^2 = 560 - 460 = 100 \quad \text{--- (2)}$$

3. \textbf{Calculate $(a - b)^2$}:
$$(a - b)^2 = 2(a^2 + b^2) - (a + b)^2 = 2(100) - (14)^2 = 200 - 196 = 4$$

Taking square root:
$$|a - b| = \sqrt{4} = 2$$

Hence, the correct option is \textbf{(B)}.
\end{tcolorbox}

\begin{tcolorbox}[solbox]
\textbf{\textcolor{primaryDark}{Question 11:}} \hfill \textbf{\textcolor{l1Green}{Correct Answer: C}}\\
\vspace{4pt}
1. \textbf{Tangent condition for the parabola} $y^2 = 4ax$ with $4a = 8 \implies a = 2$:
A tangent with slope $m$ has the equation:
$$y = mx + \frac{a}{m} = mx + \frac{2}{m} \implies mx - y + \frac{2}{m} = 0$$
Thus, the intercept is $c = \dfrac{2}{m}$.

2. \textbf{Tangent condition for the circle} $x^2 + y^2 = 2$ (radius $r = \sqrt{2}$):
The perpendicular distance from the center $(0,0)$ to the tangent line must equal the radius $r$:
$$\frac{\left|0 - 0 + \dfrac{2}{m}\right|}{\sqrt{m^2 + (-1)^2}} = \sqrt{2}$$
$$\frac{2/m}{\sqrt{m^2 + 1}} = \sqrt{2}$$

Squaring both sides:
$$\frac{4}{m^2(m^2 + 1)} = 2 \implies m^2(m^2 + 1) = 2$$
$$m^4 + m^2 - 2 = 0$$
$$(m^2 + 2)(m^2 - 1) = 0$$

Since $m \in \mathbb{R}$, $m^2 = 1$. Given $m > 0$, we have $m = 1$.

3. \textbf{Find $c$ and $m + c$}:
$$c = \frac{2}{m} = \frac{2}{1} = 2$$
$$m + c = 1 + 2 = 3$$

Hence, the correct option is \textbf{(C)}.
\end{tcolorbox}

\begin{tcolorbox}[solbox]
\textbf{\textcolor{primaryDark}{Question 12:}} \hfill \textbf{\textcolor{l1Green}{Correct Answer: B}}\\
\vspace{4pt}
Using the standard combinatorial identity:
$$\frac{^nC_r}{r + 1} = \frac{^{n+1}C_{r+1}}{n + 1}$$

Setting $n = 20$:
$$\frac{^{20}C_r}{r + 1} = \frac{^{21}C_{r+1}}{21}$$

Substituting into the summation:
$$S = \sum_{r=0}^{20} (-1)^r \frac{^{20}C_r}{r + 1} = \frac{1}{21} \sum_{r=0}^{20} (-1)^r \, ^{21}C_{r+1}$$

Let index $k = r + 1$. When $r = 0, k = 1$; when $r = 20, k = 21$.
Since $(-1)^r = -(-1)^{r+1} = -(-1)^k$:
$$S = -\frac{1}{21} \sum_{k=1}^{21} (-1)^k \, ^{21}C_k$$

From the binomial expansion of $(1 - 1)^{21}$:
$$(1 - 1)^{21} = \sum_{k=0}^{21} (-1)^k \, ^{21}C_k = ^{21}C_0 + \sum_{k=1}^{21} (-1)^k \, ^{21}C_k = 0$$
$$1 + \sum_{k=1}^{21} (-1)^k \, ^{21}C_k = 0 \implies \sum_{k=1}^{21} (-1)^k \, ^{21}C_k = -1$$

Substituting this back:
$$S = -\frac{1}{21} (-1) = \frac{1}{21}$$

Hence, the correct option is \textbf{(B)}.
\end{tcolorbox}

\begin{tcolorbox}[solbox]
\textbf{\textcolor{primaryDark}{Question 13:}} \hfill \textbf{\textcolor{l1Green}{Correct Answer: A}}\\
\vspace{4pt}
Three vectors $\vec{a}, \vec{b}, \vec{c}$ are coplanar if and only if their scalar triple product vanishes:
$$[\vec{a} \quad \vec{b} \quad \vec{c}] = 0$$

Expressing as a $3 \times 3$ determinant:
$$\begin{vmatrix} 1 & 1 & 1 \\ 1 & -1 & 2 \\ x & x - 2 & -1 \end{vmatrix} = 0$$

Expanding along row 1:
$$1 [(-1)(-1) - 2(x - 2)] - 1 [1(-1) - 2x] + 1 [1(x - 2) - (-1)x] = 0$$
$$[1 - 2x + 4] - [-1 - 2x] + [x - 2 + x] = 0$$
$$(5 - 2x) + (1 + 2x) + (2x - 2) = 0$$
$$(5 + 1 - 2) + (-2x + 2x + 2x) = 0$$
$$4 + 2x = 0 \implies 2x = -4 \implies x = -2$$

Hence, the correct option is \textbf{(A)}.
\end{tcolorbox}

\begin{tcolorbox}[solbox]
\textbf{\textcolor{primaryDark}{Question 14:}} \hfill \textbf{\textcolor{l1Green}{Correct Answer: B}}\\
\vspace{4pt}
1. \textbf{Find points of intersection}:
$$x^2 = 2x + 3 \implies x^2 - 2x - 3 = 0$$
$$(x - 3)(x + 1) = 0 \implies x = -1 \quad \text{and} \quad x = 3$$

2. \textbf{Area integral}:
On the interval $[-1, 3]$, the straight line lies above the parabola ($2x + 3 \ge x^2$):
$$A = \int_{-1}^{3} \left( (2x + 3) - x^2 \right)\, dx$$
$$A = \left[ x^2 + 3x - \frac{x^3}{3} \right]_{-1}^{3}$$

At $x = 3$:
$$\left[ 3^2 + 3(3) - \frac{3^3}{3} \right] = 9 + 9 - 9 = 9$$

At $x = -1$:
$$\left[ (-1)^2 + 3(-1) - \frac{(-1)^3}{3} \right] = 1 - 3 + \frac{1}{3} = -2 + \frac{1}{3} = -\frac{5}{3}$$

Subtracting the limits:
$$A = 9 - \left(-\frac{5}{3}\right) = 9 + \frac{5}{3} = \frac{32}{3}$$

Hence, the correct option is \textbf{(B)}.
\end{tcolorbox}

\begin{tcolorbox}[solbox]
\textbf{\textcolor{primaryDark}{Question 15:}} \hfill \textbf{\textcolor{l1Green}{Correct Answer: B}}\\
\vspace{4pt}
A number is divisible by 4 if and only if the number formed by its last two digits is divisible by 4.

1. \textbf{Valid distinct 2-digit pairs divisible by 4 from $\{1, 2, 3, 4, 5, 6, 7\}$}:
- Tens digit 1: $12, 16$ (2 pairs)
- Tens digit 2: $24$ (1 pair)
- Tens digit 3: $32, 36$ (2 pairs)
- Tens digit 4: none ($44$ repeats 4)
- Tens digit 5: $52, 56$ (2 pairs)
- Tens digit 6: $64$ (1 pair)
- Tens digit 7: $72, 76$ (2 pairs)
Total valid ending pairs: $2 + 1 + 2 + 0 + 2 + 1 + 2 = 10$ pairs.

2. \textbf{Filling the first 3 positions}:
For each valid 2-digit ending, 2 distinct digits have been chosen from 7.
There are $7 - 2 = 5$ remaining digits to arrange into the first 3 positions:
$$^5P_3 = 5 \times 4 \times 3 = 60 \text{ ways}$$

3. \textbf{Total count}:
$$N = 10 \times 60 = 600$$

Hence, the correct option is \textbf{(B)}.
\end{tcolorbox}

\begin{tcolorbox}[solbox]
\textbf{\textcolor{primaryDark}{Question 16:}} \hfill \textbf{\textcolor{l1Green}{Correct Answer: B}}\\
\vspace{4pt}
For an ellipse $\dfrac{x^2}{a^2} + \dfrac{y^2}{b^2} = 1$, any point in the first quadrant has coordinates $(a\cos\theta, b\sin\theta)$ with $\theta \in \left(0, \dfrac{\pi}{2}\right)$.

Here $a = 4$ and $b = 3$.

A rectangle with sides parallel to the axes has vertices $(\pm 4\cos\theta, \pm 3\sin\theta)$.
- Length along $x$-axis $= 2(4\cos\theta) = 8\cos\theta$
- Width along $y$-axis $= 2(3\sin\theta) = 6\sin\theta$

The area function is:
$$A(\theta) = (8\cos\theta)(6\sin\theta) = 48\cos\theta\sin\theta = 24\sin(2\theta)$$

Since $-1 \le \sin(2\theta) \le 1$, the maximum value of $\sin(2\theta)$ is 1 (attained at $\theta = \dfrac{\pi}{4}$).

Therefore, the maximum area is:
$$A_{\max} = 24(1) = 24$$

(In general, for any ellipse with semi-axes $a$ and $b$, $A_{\max} = 2ab = 2(4)(3) = 24$).

Hence, the correct option is \textbf{(B)}.
\end{tcolorbox}

\begin{tcolorbox}[solbox]
\textbf{\textcolor{primaryDark}{Question 17:}} \hfill \textbf{\textcolor{l1Green}{Correct Answer: A}}\\
\vspace{4pt}
1. \textbf{Continuity at $x = 0$}:
For $f(x)$ to be continuous at $x = 0$:
$$\lim_{x \to 0} f(x) = f(0)$$
$$\lim_{x \to 0} \left[ x^2 \sin\left(\frac{1}{x}\right) + ax \right] = b$$
Since $\left|\sin\left(\frac{1}{x}\right)\right| \le 1$:
$$\lim_{x \to 0} x^2 \sin\left(\frac{1}{x}\right) = 0 \quad \text{and} \quad \lim_{x \to 0} ax = 0$$
Therefore, $0 = b \implies b = 0$.

2. \textbf{Differentiability at $x = 0$}:
By the definition of derivative:
$$f'(0) = \lim_{h \to 0} \frac{f(h) - f(0)}{h} = \lim_{h \to 0} \frac{h^2 \sin\left(\dfrac{1}{h}\right) + ah - 0}{h}$$
$$f'(0) = \lim_{h \to 0} \left[ h \sin\left(\frac{1}{h}\right) + a \right]$$
Since $\lim_{h \to 0} h \sin(1/h) = 0$:
$$f'(0) = 0 + a = a$$

Given $f'(0) = 3$, we have $a = 3$.

Thus, the ordered pair is $(a, b) = (3, 0)$.

Hence, the correct option is \textbf{(A)}.
\end{tcolorbox}

\begin{tcolorbox}[solbox]
\textbf{\textcolor{primaryDark}{Question 18:}} \hfill \textbf{\textcolor{l1Green}{Correct Answer: A}}\\
\vspace{4pt}
Express the argument of the inverse tangent function in difference form:
$$\frac{1}{r^2 + r + 1} = \frac{(r + 1) - r}{1 + r(r + 1)}$$

Using the identity $\tan^{-1}\left(\dfrac{x - y}{1 + xy}\right) = \tan^{-1}x - \tan^{-1}y$ for positive $x, y$:
$$\tan^{-1}\left(\frac{1}{r^2 + r + 1}\right) = \tan^{-1}(r + 1) - \tan^{-1}(r)$$

Summing from $r = 1$ to $10$:
$$S = \sum_{r=1}^{10} \left[ \tan^{-1}(r + 1) - \tan^{-1}(r) \right]$$
$$S = (\tan^{-1}2 - \tan^{-1}1) + (\tan^{-1}3 - \tan^{-1}2) + \dots + (\tan^{-1}11 - \tan^{-1}10)$$
$$S = \tan^{-1}(11) - \tan^{-1}(1)$$

Now compute $\tan(S)$:
$$\tan(S) = \tan\left(\tan^{-1}(11) - \tan^{-1}(1)\right) = \frac{11 - 1}{1 + 11(1)} = \frac{10}{1 + 11} = \frac{10}{12} = \frac{5}{6}$$

Hence, the correct option is \textbf{(A)}.
\end{tcolorbox}

\begin{tcolorbox}[solbox]
\textbf{\textcolor{primaryDark}{Question 19:}} \hfill \textbf{\textcolor{l1Green}{Correct Answer: B}}\\
\vspace{4pt}
1. \textbf{Analyze both circles}:
- Circle $C_1$: $x^2 + y^2 + 4x + 6y - 12 = 0$
  - Center: $(-2, -3)$
  - Radius $R_1 = \sqrt{(-2)^2 + (-3)^2 - (-12)} = \sqrt{4 + 9 + 12} = \sqrt{25} = 5$

- Circle $C_2$: $x^2 + y^2 + 4x + 6y - 4 = 0$
  - Center: $(-2, -3)$
  - Radius $R_2 = \sqrt{(-2)^2 + (-3)^2 - (-4)} = \sqrt{4 + 9 + 4} = \sqrt{17}$

Both circles share the same center $(-2, -3)$ and are concentric.

2. \textbf{Tangent Length from point on $C_1$ to $C_2$}:
For any point $P$ on $C_1$, the distance from $P$ to the common center is $d = R_1 = 5$.

The length of the tangent $L$ from point $P$ to circle $C_2$ forms a right triangle with radius $R_2$:
$$L = \sqrt{d^2 - R_2^2} = \sqrt{R_1^2 - R_2^2}$$
$$L = \sqrt{5^2 - (\sqrt{17})^2} = \sqrt{25 - 17} = \sqrt{8} = 2\sqrt{2}$$

Hence, the correct option is \textbf{(B)}.
\end{tcolorbox}

\begin{tcolorbox}[solbox]
\textbf{\textcolor{primaryDark}{Question 20:}} \hfill \textbf{\textcolor{l1Green}{Correct Answer: A}}\\
\vspace{4pt}
Let $B_1$ be the event of selecting Box I, $B_2$ be the event of selecting Box II, and $R$ be the event that a red ball is drawn.

Given:
$$P(B_1) = P(B_2) = \frac{1}{2}$$
$$P(R | B_1) = \frac{4}{4 + 6} = \frac{4}{10} = \frac{2}{5}$$
$$P(R | B_2) = \frac{5}{5 + 5} = \frac{5}{10} = \frac{1}{2}$$

Applying Bayes' Theorem:
$$P(B_1 | R) = \frac{P(B_1) P(R | B_1)}{P(B_1) P(R | B_1) + P(B_2) P(R | B_2)}$$
$$P(B_1 | R) = \frac{\dfrac{1}{2} \times \dfrac{4}{10}}{\dfrac{1}{2} \times \dfrac{4}{10} + \dfrac{1}{2} \times \dfrac{5}{10}} = \frac{\dfrac{4}{10}}{\dfrac{4}{10} + \dfrac{5}{10}} = \frac{4}{4 + 5} = \frac{4}{9}$$

Hence, the correct option is \textbf{(A)}.
\end{tcolorbox}

\begin{tcolorbox}[solbox]
\textbf{\textcolor{primaryDark}{Question 21:}} \hfill \textbf{\textcolor{l1Green}{Correct Answer: 12}}\\
\vspace{4pt}
Given the equation:
$$\int_{0}^{x^2} f(t)\, dt = \frac{2}{5}x^5 + \frac{4}{3}x^3$$

Differentiating both sides with respect to $x$ using the Leibniz Rule for differentiating under the integral sign:
$$\frac{d}{dx}\left[ \int_{0}^{x^2} f(t)\, dt \right] = f(x^2) \cdot \frac{d}{dx}(x^2) - f(0) \cdot \frac{d}{dx}(0) = f(x^2) \cdot 2x$$

Taking derivative of the right-hand side:
$$\frac{d}{dx}\left[ \frac{2}{5}x^5 + \frac{4}{3}x^3 \right] = 2x^4 + 4x^2$$

Equating both derivatives:
$$2x \cdot f(x^2) = 2x^4 + 4x^2$$

Since $x > 0$, dividing both sides by $2x$:
$$f(x^2) = x^3 + 2x$$

We need to find $f(4)$. Setting $x^2 = 4 \implies x = 2$ (since $x > 0$):
$$f(4) = (2)^3 + 2(2) = 8 + 4 = 12$$

Hence, the correct integer answer is \textbf{12}.
\end{tcolorbox}

\begin{tcolorbox}[solbox]
\textbf{\textcolor{primaryDark}{Question 22:}} \hfill \textbf{\textcolor{l1Green}{Correct Answer: 20}}\\
\vspace{4pt}
Let matrix $A$ be decomposed into $A = I + N$, where:
$$I = \begin{bmatrix} 1 & 0 \\ 0 & 1 \end{bmatrix}, \quad N = \begin{bmatrix} 0 & 2 \\ 0 & 0 \end{bmatrix}$$

Notice that $N^2 = \begin{bmatrix} 0 & 2 \\ 0 & 0 \end{bmatrix}\begin{bmatrix} 0 & 2 \\ 0 & 0 \end{bmatrix} = \begin{bmatrix} 0 & 0 \\ 0 & 0 \end{bmatrix} = O$.

Since $I$ and $N$ commute ($IN = NI = N$), we can apply the Binomial Theorem:
$$A^{10} = (I + N)^{10} = I^{10} + \binom{10}{1} I^9 N + \binom{10}{2} I^8 N^2 + \dots + N^{10}$$
$$A^{10} = I + 10N + O = \begin{bmatrix} 1 & 0 \\ 0 & 1 \end{bmatrix} + 10 \begin{bmatrix} 0 & 2 \\ 0 & 0 \end{bmatrix} = \begin{bmatrix} 1 & 20 \\ 0 & 1 \end{bmatrix}$$

Comparing with $A^{10} = \begin{bmatrix} 1 & k \\ 0 & 1 \end{bmatrix}$, we obtain:
$$k = 20$$

Hence, the correct integer answer is \textbf{20}.
\end{tcolorbox}

\begin{tcolorbox}[solbox]
\textbf{\textcolor{primaryDark}{Question 23:}} \hfill \textbf{\textcolor{l1Green}{Correct Answer: 3}}\\
\vspace{4pt}
1. \textbf{Determine $|\vec{a}|$ and $|\vec{c}|$}:
$$|\vec{a}|^2 = 2^2 + 1^2 + (-2)^2 = 4 + 1 + 4 = 9 \implies |\vec{a}| = 3$$

Given $|\vec{c} - \vec{a}| = 2\sqrt{2}$:
$$|\vec{c} - \vec{a}|^2 = |\vec{c}|^2 + |\vec{a}|^2 - 2(\vec{a} \cdot \vec{c}) = (2\sqrt{2})^2 = 8$$

Substitute $|\vec{a}|^2 = 9$ and $\vec{a} \cdot \vec{c} = |\vec{c}|$:
$$|\vec{c}|^2 + 9 - 2|\vec{c}| = 8$$
$$|\vec{c}|^2 - 2|\vec{c}| + 1 = 0 \implies (|\vec{c}| - 1)^2 = 0 \implies |\vec{c}| = 1$$

2. \textbf{Compute $\vec{a} \times \vec{b}$}:
$$\vec{a} \times \vec{b} = \begin{vmatrix} \hat{i} & \hat{j} & \hat{k} \\ 2 & 1 & -2 \\ 1 & 1 & 0 \end{vmatrix}$$
$$= \hat{i}(0 - (-2)) - \hat{j}(0 - (-2)) + \hat{k}(2 - 1) = 2\hat{i} - 2\hat{j} + \hat{k}$$
$$|\vec{a} \times \vec{b}| = \sqrt{2^2 + (-2)^2 + 1^2} = \sqrt{4 + 4 + 1} = \sqrt{9} = 3$$

3. \textbf{Compute magnitude of the cross product}:
Since $\vec{c}$ is perpendicular to $\vec{a} \times \vec{b}$, the angle between them is $\theta = 90^\circ$:
$$|(\vec{a} \times \vec{b}) \times \vec{c}| = |\vec{a} \times \vec{b}| \cdot |\vec{c}| \cdot \sin(90^\circ) = 3 \times 1 \times 1 = 3$$

Hence, the correct integer answer is \textbf{3}.
\end{tcolorbox}

\begin{tcolorbox}[solbox]
\textbf{\textcolor{primaryDark}{Question 24:}} \hfill \textbf{\textcolor{l1Green}{Correct Answer: 14}}\\
\vspace{4pt}
For the standard parabola $y^2 = 4ax$ with $a = 1$:

1. \textbf{Focal Distance Property}:
The distance of any point $(x, y)$ on the parabola from the focus is given by:
$$SP = x + a = x + 1$$

For endpoints $P(x_1, y_1)$ and $Q(x_2, y_2)$ of a focal chord:
$$PQ = SP + SQ = (x_1 + a) + (x_2 + a) = x_1 + x_2 + 2a$$

2. \textbf{Substitute $a = 1$ and $PQ = 16$}:
$$16 = x_1 + x_2 + 2(1)$$
$$x_1 + x_2 + 2 = 16 \implies x_1 + x_2 = 14$$

(Alternatively, using parameter $t$: $x_1 = t^2, x_2 = 1/t^2$, length $PQ = (t + 1/t)^2 = t^2 + 1/t^2 + 2 = 16 \implies t^2 + 1/t^2 = 14$).

Hence, the correct integer answer is \textbf{14}.
\end{tcolorbox}

\begin{tcolorbox}[solbox]
\textbf{\textcolor{primaryDark}{Question 25:}} \hfill \textbf{\textcolor{l1Green}{Correct Answer: 15}}\\
\vspace{4pt}
Let the three roots of the cubic equation be $(a - d)$, $a$, and $(a + d)$, with common difference $d > 0$.

1. \textbf{Sum of roots (Vieta's formula)}:
$$(a - d) + a + (a + d) = -(-9) = 9$$
$$3a = 9 \implies a = 3$$

2. \textbf{Sum of pairwise products of roots}:
$$(a - d)a + a(a + d) + (a - d)(a + d) = 23$$
$$(3 - d)(3) + 3(3 + d) + (3^2 - d^2) = 23$$
$$(9 - 3d) + (9 + 3d) + (9 - d^2) = 23$$
$$27 - d^2 = 23 \implies d^2 = 4 \implies d = 2$$

3. \textbf{Roots of the cubic}:
The roots are:
$$x_1 = a - d = 3 - 2 = 1$$
$$x_2 = a = 3$$
$$x_3 = a + d = 3 + 2 = 5$$
These roots $1, 3, 5$ are distinct positive integers in arithmetic progression.

4. \textbf{Constant term $c$}:
By Vieta's formula for the product of roots:
$$\text{Product of roots} = -(-c) = c$$
$$c = (1)(3)(5) = 15$$

Hence, the correct integer answer is \textbf{15}.
\end{tcolorbox}

\end{document}